\section{R-DIAG-001: a free cyclic action does not give three-integral local terms}

This elementary weighted hypersurface records the failure of one
proposed shortcut at the remaining pole gate. It is not the tensor
quotient, and it is not a counterexample to the modular pole-transfer
theorem. No novelty claim is made for the underlying toric hypersurface.

Let $M\ge2$ be a power of two and put
\[
 R=\mathbb C[x,y,z,w]/(xyz-w^2),\qquad
 \deg x=\deg y=\deg z=3M,\qquad \deg w=9M/2.
                                                        \tag{R-DIAG-001a}
\]
Let $\tau$ cycle $x,y,z$ and fix $w$. The scheme
$X=\operatorname{Spec}R$ is irreducible and normal. Its grading-fixed
schemes satisfy
\[
 X^{\mu_M}=X^{\mu_{3M}}=\{w=0,\ xyz=0\}.
                                                        \tag{R-DIAG-001b}
\]
They are the three reduced coordinate planes. Their generic points
are smooth points of $X$, and $\tau$ cycles the components with no
nonzero fixed point on their union. Nevertheless, at every primitive
$M$th root, with local parameter $u=1-t^M$, the Hilbert series has
\[
 H_R(t)=\frac1{6u^2}+\frac5{24u}+O(1).
                                                        \tag{R-DIAG-001c}
\]
Both coefficients have three-adic valuation $-1$.

\begin{proof}
Let $F\subset\{\pm1\}^3$ be the subgroup whose coordinate product
is one. A monomial in $\mathbb C[a,b,c]$ is $F$-invariant exactly
when its three exponents have the same parity. Consequently
\[
 \mathbb C[a,b,c]^F
 =\mathbb C[a^2,b^2,c^2]\oplus abc\,\mathbb C[a^2,b^2,c^2]
 \simeq R.
\]
The even and odd monomial parts are disjoint, and multiplication
gives $(abc)^2=a^2b^2c^2$, proving the stated presentation. This is
a domain. It is normal because any fraction in its fraction field
integral over it is integral over the normal polynomial ring,
therefore belongs to that ring, and remains $F$-invariant. This
also proves normality directly for the particular presentation.

A primitive grading element of order $M$ acts as $+1$ on $x,y,z$
and $-1$ on $w$; the same is true for order $3M$. This proves
(R-DIAG-001b). The squarefree equation $xyz=0$ defines the reduced
union of the three planes. A fixed point of their cyclic permutation
has $x=y=z$, so $x^3=0$ and it is the origin. On the plane $z=0$,
the open subset $xy\ne0$ is smooth in $X$ because the derivative
of the defining equation with respect to $z$ is $xy$; the other
planes are the same by symmetry.

The monic relation in $w$ makes $R$ a free graded module with basis
$1,w$ over $\mathbb C[x,y,z]$. Hence
\[
 H_R(t)=\frac{1+t^{9M/2}}{(1-t^{3M})^3}.
                                                        \tag{R-DIAG-001d}
\]
Near a primitive $M$th root, $t^{M/2}=-(1-u)^{1/2}$, with the
analytic branch equal to one at $u=0$. Thus the right side is
\[
 \frac{1-(1-u)^{9/2}}{\bigl(1-(1-u)^3\bigr)^3}.
\]
Its numerator is $\tfrac92u-\tfrac{63}8u^2+O(u^3)$, and the
denominator is $27u^3(1-3u+O(u^2))$. Dividing gives
(R-DIAG-001c). At a generic point of each fixed plane, the moving
tangent coordinate has eigenvalue $-1$, so its inverse Euler
factor is $1/2$. Each plane has coefficient $1/9$ in the
normalization $(1-t^M)^{-2}$. The three contributions therefore
sum to $3(1/2)(1/9)=1/6$, agreeing with the direct series.
\end{proof}

For completeness, the same two-term module decomposition computes
the cyclic trace exactly:
\[
 \sum_d\operatorname{tr}(\tau\mid R_d)t^d
 =\frac{1+t^{9M/2}}{1-t^{9M}}
 =\frac1{1-t^{9M/2}}.
                                                        \tag{R-DIAG-001e}
\]
This trace has no pole at an order $M3^j$. In particular the
example does not satisfy the auxiliary-pole hypothesis of
PNEW-013. Its precise purpose is different: absence of nonzero
cyclic fixed points, even on generically smooth fixed components
of a normal irreducible hypersurface, does not make their full
Laurent principal part three-integral. Additional tensor-specific
information is required to remove K-031's collision exception.

The exact rational identities and the displayed Laurent coefficients
are replayed by
\nolinkurl{branches/root/verify_cyclic_free_collision.py}.
The proof, rather than its finite character sample, establishes
the all-degree trace identity and the geometry.
